\section{Background}
\label{sec:background}

\subsection{OSEK/VDX Task, Scheduler, and Event Semantics}

The seed introduction material is centered on an observation that still holds for the current
branch: OSEK applications are difficult to verify because task scheduling is explicit and API
driven rather than an afterthought. In OSEK/VDX, tasks have well-defined states, priorities,
and scheduling policies, and the operating system mediates transitions between those states
through standardized services~\cite{osek_spec}. Basic tasks move among suspended, ready, and
running states. Extended tasks add waiting behavior and event masks, which means APIs such as
\texttt{SetEvent}, \texttt{GetEvent}, \texttt{ClearEvent}, and \texttt{WaitEvent} modify both
control flow and visible state.

This matters because the scheduler is not nondeterministic in the same way as a generic
threading runtime. OSEK code frequently relies on explicit priorities, ready queues, and
preemption rules. The benchmark cases in this repository also use OIL metadata to define task
identities and priorities, so a verification flow that ignores the OIL configuration loses a
large fraction of the execution model before symbolic execution even begins.

\subsection{Baseline CBMC Architecture}

CBMC follows a standard three-stage structure~\cite{cbmc_architecture}. First, a front-end
and \texttt{goto-cc} convert source files into a GOTO program. Second,
\texttt{goto-symex} symbolically executes that model and emits an SSA-style equation. Third,
the solver backend discharges the equation and optionally reconstructs a counterexample trace.
This is a strong baseline for ordinary C verification because it already understands low-level
control flow, bit-precise data, and solver-backed safety checks.

\begin{figure}[t]
  \centering
  \begin{tikzpicture}[
  node distance=7mm,
  box/.style={
    draw,
    rounded corners=2pt,
    align=center,
    minimum width=34mm,
    minimum height=10mm,
    fill=gray!8
  },
  arrow/.style={-{Latex[length=2mm]},thick}
]
  \node[box] (src) {C sources\\and harness};
  \node[box,below=of src] (goto) {Front-end\\and \texttt{goto-cc}};
  \node[box,below=of goto] (model) {GOTO model};
  \node[box,below=of model] (symex) {\texttt{goto-symex}\\SSA equation};
  \node[box,below=of symex] (solver) {SAT/SMT\\solver};
  \node[box,below=of solver] (trace) {Result and\\counterexample};

  \draw[arrow] (src) -- (goto);
  \draw[arrow] (goto) -- (model);
  \draw[arrow] (model) -- (symex);
  \draw[arrow] (symex) -- (solver);
  \draw[arrow] (solver) -- (trace);
\end{tikzpicture}

  \caption{Baseline CBMC flow before the OSEK-aware and manifest-driven extensions described
  in this report.}
  \label{fig:baseline-cbmc}
\end{figure}

The limitation is not that CBMC lacks symbolic execution, but that a stock flow does not
automatically know which API calls are scheduler actions and which source statements are worth
instrumenting for ISR interference. Those are domain-specific choices. In this branch, they
are made explicitly instead of being hidden in monolithic harness code.

\subsection{Three Families of Interrupt-Aware Verification}

The background PDF reviews several styles of model checking for concurrent software. That
survey still gives a useful way to position the current implementation, provided the
terminology is updated to match the code in this repository.

\paragraph{Explicit-state exploration.}
Classical explicit-state tools such as SPIN build and traverse an explicit transition system,
which makes them a natural reference point for state-based concurrency reasoning~\cite{spin_site}.
Their strength is transparent control-state exploration; their weakness for this project is
that the OSEK scheduler and ISR placement policy still have to be modeled carefully to avoid
spurious interleavings.

\paragraph{SMT- and source-transformation-oriented flows.}
The benchmark corpora in this repository include artifacts from the CPROVER interrupt work and
from IntAbs~\cite{icbmc_site,intabs_repo}. Those families combine symbolic reasoning with
specialized source transformation or interrupt abstraction so that interleavings can be
restricted to relevant program points. The improved pipeline in this branch is closest to this
family, but it differs in two important ways: the OSEK scheduler semantics are integrated
inside CBMC rather than being left entirely to source transformation, and the static analysis
phase is cleanly separated from rewriting by a manifest boundary.

\paragraph{Sequentialization.}
Another well-known family translates concurrent behavior into a sequential model that can be
fed to an existing verifier. The background material emphasizes this as a scalability tactic
for branch-heavy applications. The current branch does not sequentialize the whole OSEK
system. Instead, it preserves CBMC's standard symbolic-execution core and inserts only the
scheduler and ISR modeling that are needed for the chosen project configuration.

These families motivate the design choices in \Cref{sec:method}. The branch documented here is
best understood as a hybrid engineering solution: keep CBMC's native GOTO and symex pipeline,
add OSEK semantics exactly where scheduler actions matter, and limit ISR modeling with a
static manifest rather than with blanket asynchronous injection.
